\section{训练流程}

\subsection{GRPO 风格的循环}

脚本中的训练循环可以概括为：
\begin{enumerate}
    \item 采样一批 prompts；
    \item 对每个 prompt 采样多个 responses；
    \item 计算 rewards 和 deltas；
    \item 可选地计算 reference model 的 log probs；
\item 用当前 policy 的 log probs 算 loss；
\item 反向传播，更新参数。
\end{enumerate}


\subsection{实验现象}

课程脚本尝试了三种设置：
\begin{itemize}
    \item raw rewards；
    \item centered rewards；
    \item normalized rewards。
\end{itemize}

结论是：
\begin{itemize}
    \item raw rewards 往往学得慢，容易卡住；
    \item centered rewards 更稳定一些；
    \item normalized rewards 不一定有明显增益，甚至可能引入不必要的偏差；
    \item 这类 toy RL 任务本身就很容易卡在局部最优。
\end{itemize}

\begin{table}[H]
\centering
\begin{tabular}{@{}p{3.2cm}p{4.6cm}p{4.8cm}@{}}
\toprule
\textbf{设置} & \textbf{观察} & \textbf{含义} \\
\midrule
raw rewards & 更新噪声大，训练容易停滞 & 绝对 reward 直接进梯度，方差过高 \\
centered rewards & 学习更稳定，负样本也有作用 & 相对比较比绝对分数更有信息 \\
normalized rewards & 有时接近 centered 的效果 & 标准差缩放主要是稳定尺度 \\
\bottomrule
\end{tabular}
\caption{三种实验设置的经验结论：GRPO 这类方法的关键通常是“先把信号做稳”}
\end{table}


\begin{warningbox}{RL 不是“设个 reward 就会自动学会”}
即便任务和 reward 都写得很清楚，policy gradient 仍然可能因为方差、初始化和超参问题而学不动。
\end{warningbox}

